\documentclass[12pt]{article}
\usepackage[margin=1in]{geometry}
\usepackage{setspace}
\onehalfspacing{}
\usepackage[dvipsnames,table,xcdraw]{xcolor} % colors

% Start of preamble
%==========================================================================================%
% Required to support mathematical unicode
\usepackage[warnunknown, fasterrors, mathletters]{ucs}
\usepackage[utf8x]{inputenc}
\usepackage{R:/sty/quiver}

% Standard mathematical typesetting packages
\usepackage{amsmath,amssymb,amscd,amsthm,amsxtra}
\usepackage{mathtools,mathrsfs,xparse, pxfonts}

% Symbol and utility packages
\usepackage{cancel, textcomp}
\usepackage[mathscr]{euscript}
\usepackage[nointegrals]{wasysym}
\usepackage{apacite}

% Extras
\usepackage{physics}  % Lots of useful shortcuts and macros
\usepackage{tikz-cd}  % For drawing commutative diagrams easily
\usepackage{microtype}  % Minature font tweaks
%\usepackage{pgfplots} % plots

\usepackage{enumitem}
\usepackage{titling}

\usepackage{graphicx}

%\usepackage{quiver}

% Fancy theorems due to @intuitively on discord
\usepackage{mdframed}
\newmdtheoremenv[
backgroundcolor=NavyBlue!30,
linewidth=2pt,
linecolor=NavyBlue,
topline=false,
bottomline=false,
rightline=false,
innertopmargin=10pt,
innerbottommargin=10pt,
innerrightmargin=10pt,
innerleftmargin=10pt,
skipabove=\baselineskip,
skipbelow=\baselineskip]{mytheorem}{Theorem}

\newenvironment{theorem}{\begin{mytheorem}}{\end{mytheorem}}

\newtheorem{corollary}{Corollary}
\newtheorem{lemma}{Lemma}

\newtheoremstyle{definitionstyle}
{\topsep}%
{\topsep}%
{}%
{}%
{\bfseries}%
{.}%
{.5em}%
{}%
\theoremstyle{definitionstyle}
\newmdtheoremenv[
backgroundcolor=Violet!30,
linewidth=2pt,
linecolor=Violet,
topline=false,
bottomline=false,
rightline=false,
innertopmargin=10pt,
innerbottommargin=10pt,
innerrightmargin=10pt,
innerleftmargin=10pt,
skipabove=\baselineskip,
skipbelow=\baselineskip,
]{mydef}{Definition}
\newenvironment{definition}{\begin{mydef}}{\end{mydef}}

\newtheorem*{remark}{Remark}

\newtheorem*{example}{Example}
\newtheorem*{claim}{Claim}

% Common shortcuts
\def\mbb#1{\mathbb{#1}}
\def\mfk#1{\mathfrak{#1}}

\def\bN{\mbb{N}}
\def\C{\mbb{C}}
\def\R{\mbb{R}}
\def\bQ{\mbb{Q}}
\def\bZ{\mbb{Z}}
\def\cph{\varphi}
\renewcommand{\th}{\theta}
\def\ve{\varepsilon}
\newcommand{\mg}[1]{\| #1 \|}

% Often helpful macros
\newcommand{\floor}[1]{\left\lfloor#1\right\rfloor}
\newcommand{\ceil}[1]{\left\lceil#1\right\rceil}
\renewcommand{\qed}{\hfill\qedsymbol}
\renewcommand{\ip}[1]{\langle#1\rangle}
\newcommand{\seq}[2]{\qty(#1_#2)_{#2=1}^{\infty}}

\newcommand{\SET}[1]{\Set{\mskip-\medmuskip #1 \mskip-\medmuskip}}

% End of preamble
%==========================================================================================%

% Start of commands specific to this file
%==========================================================================================%

\usepackage{braket}
\newcommand{\Z}{\mbb Z}
\newcommand{\gen}[1]{\left\langle #1 \right\rangle}
\newcommand{\nsg}{\trianglelefteq}
\newcommand{\F}{\mbb F}
\newcommand{\Aut}{\mathrm{Aut}}
\newcommand{\sepdeg}[1]{[#1]_{\mathrm{sep}}}
\newcommand{\Q}{\mbb Q}
\newcommand{\Gal}{\mathrm{Gal}\qty}
\newcommand{\id}{\mathrm{id}}
\newcommand{\Hom}{\mathrm{Hom}_R}
\renewcommand{\S}{\mbb S}

%==========================================================================================%
% End of commands specific to this file

\title{Math Template}
\date{\today}
\author{Rohan Mukherjee}

\begin{document}
    \maketitle
    \subsection*{Problem 1.}
    \begin{theorem}
        Let $G$ be a finite graph and define $d$ to be the smallest integer such that $G$ has a unit distance representation in $\R^d$. Then $\log(\chi(G))/2 \leq d \leq 2\chi(G)$. 
    \end{theorem}
    We shall show that if the points are required to be distinct, then $d \leq 2\chi(G)$. Let $G$ be a graph with chromatic number $\chi(G)$ and let $V_1, \ldots, V_k$ be the partition of $V$ into $\chi(G)$ parts where each of the $V_i$ are independent. We embed $G$ into $\R^{2\chi(G)}$ by doing the following: for each $v_j \in V_i$, we assign a coordinate to $v$ as $(0, 0, \ldots, x_{i,j}, y_{i,j}, 0, 0, \ldots, 0, 0)$ where $x_{i,j}, y_{i,j}$ are in the $2i$th and $2i+1$th spot respectively, satisfying $x_{i,j}^2 + y_{i,j}^2 = 1/2$. We can do this since the final equation clearly has infinitely many solutions and our grpah is finite. Then we see that if $u \to v$ is an edge with $u \in V_k$ and $v \in V_\ell$, then the distance between $u$ and $v$ in the embedding is just clearly 1. The lower bound remains the same, since we have restricted the points to be distinct, so the minimal dimension can only go up. The same subgraph of $H_d$ with chromatic number $1.2^d$ shows that this bound is tight up to constants, the same way it was used to show that the dimension of the orthogonal representation of a graph is at most $2\chi(G)$. We make this a unit distance representation by making the sphere in dimension $d$ have radius $1/2$ instead of 1. Recall that there is an edge between $u, v \in \S^{d-1}$ iff $u,v$ are orthogonal. Then we see that the distance between $u$ and $v$ is just $2$, which we made a unit distance representation by normalizing.
    
    \subsection*{Problem 2.}
    Find the dual of the following SDP:
    \begin{align*}
        \text{minimize } & x \\
        \text{subject to } & \begin{pmatrix}
            x & 1 \\ 
            1 & y
        \end{pmatrix} \succeq 0
    \end{align*}
    We notice that 
    \begin{align*}
        \begin{pmatrix}
            x & 1 \\ 
            1 & y
        \end{pmatrix} = x\begin{pmatrix}
            1 & 0 \\ 
            0 & 0
        \end{pmatrix} + y\begin{pmatrix}
            0 & 0 \\ 
            0 & 1
        \end{pmatrix} + \begin{pmatrix}
            0 & 1 \\ 
            1 & 0
        \end{pmatrix}
    \end{align*}
    Thus the dual becomes:
    \begin{align*}
        \text{minimize } & \begin{pmatrix}
            0 & 1 \\ 
            1 & 0
        \end{pmatrix} \bullet X \\
        \text{subject to } & \begin{pmatrix}
            1 & 0 \\ 
            0 & 0
        \end{pmatrix} \bullet X = 1 \\
        & \begin{pmatrix}
            0 & 0 \\ 
            0 & 1
        \end{pmatrix} \bullet X = 0 \\
        & X \succeq 0
    \end{align*}
    Write $X = \begin{pmatrix}
        a & b \\ 
        b & c
    \end{pmatrix}$. This SDP has a feasible solution: for example, taking $X = \begin{pmatrix}
        1 & a \\ 
        a & 0
    \end{pmatrix}$ satisfies the constraints and gives a value of $2a$. The above matrix is easily seen to have eigenvalues of $\frac12 \pm \frac12 \sqrt{1-4a^2}$, so taking $a^2 < 1/4$ equivalently $|a| < 1/2$ will yield a strictly feasible solution to the above SDP. Thus by strong duality the optimal value of the primal and dual are equal, and this value is easily seen to be $-1$ by our classification of the above matrix.

    \subsection*{Problem 3.}
    Let $A_1, \ldots, A_m$ be symmetric $m \times m$ matrices. Then exactly one of the following holds: (i) There are $x_1, \ldots, x_n$ so that $x_1A_1 + \cdots + x+nA_n \succeq 0$ and $\sum x_iA_i \neq 0$, (ii) There exists a matrix $Y \succ 0$ such that $A_i \cdot Y = 0$ for all $i$.

    We seek to apply Farkas lemma. Let $L = \mathrm{span}\SET{A_1, \ldots, A_n}$. By Farkas lemma, there is either a matrix $X \in L^\perp$ with $X \succ 0$, or there is a matrix $Y \in (L^\perp)^\perp = L$ so that $Y \neq 0$, and $Y \succ 0$. The first condition obviously says that there is a matrix $Y$ with $A_i \cdot Y = 0$ for all $i$. The second says that there is a matrix $Y \in L$ so that $y \succeq 0$ and $Y \neq 0$. By definition of $L$, $Y = \sum x_iA_i$ for some coefficients $x_i$, which completes the proof.

    \subsection*{Problem 4.}
    Let $A_1, \ldots, A_m$ be symmetric $m \times m$ matrices and $b_1, \ldots, b_n \in \R$. Then exactly one of the following alternatives holds: (i) There exist $x_1, \ldots, x_n \in \R$ so that $x_1A_1 + \cdots + x_nA_n \succeq 0$ and $\sum_i x_ib_i \leq 0$, and if $\sum_i x_ib_i = 0$, then $\sum_i |b_i| > 0$ and $y^TA_iy = 0$ for all $y \in \R^m$ with $(x_1A_1 + \cdots + x_nA_n)y = 0$. (ii) There exists a vector $y \in \R^m$ such that $A_i \cdot y = b_i$ for all $i$.

    The version of Farkas lemma proved in the book is the following:
    Let $\mathcal A$ be defined by the set of equations $A_i \cdot X = c_i$ for $i = 1,\ldots,n$. Then define
    \begin{align*}
        \widehat{\mathcal A} &= \SET{X \in S \; | \; A_i \cdot X = 0}, \\
        \widehat{\mathcal A^*} &= \SET{\sum x_iA_i \; | \; \sum x_ic_i = 0}, \\
        \mathcal A^* &= \SET{\sum x_iA_i \; | \; \sum x_ic_i = -1}
    \end{align*}
    Then either $\mathcal A$ or $\mathcal A^*$ contains a nonzero PSD matrix, or both $\widehat{\mathcal A}$ and $\widehat{\mathcal A^*}$ contain nonzero PSD matrices.

    Suppose that $\mathcal{A} = \SET{X \in S \; | \; A_i \cdot X = b_i}$ is so that all the $b_i$ are 0. Then as we can take $Y=0$ to show that $A_i \cdot Y = 0$ for all $i$, we have to show that the first alternative cannot hold. Suppose there were $x_1, \ldots, x_n \in \R$ so that $\sum x_i A_i \succeq 0$, $\sum x_ib_i \leq 0$, and if $\sum x_ib_i = 0$ then $\sum_i |b_i| > 0$ and $y^TAy = 0$ for all $y \in \ker \sum x_iA_i$. This follows because if we take all the $x_i = 0$, then we do have $\sum x_iA_i \succeq 0$, $\sum x_ib_i = 0 \leq 0$, while $\sum |b_i| = 0$ is not greater than 0. 
    
    So we can assume that not all of the $b_i$ are 0. Suppose that both conditions hold at the same time--i.e. that there is some $Y \succeq 0$ so that $A_i \cdot Y = b_i$ for all $i$, and that there exists $x_1, \ldots, x_n \in \R$ so that $\sum x_iA_i \succeq 0$, $\sum_i x_ib_i \leq 0$, and if $\sum_i x_ib_i = 0$ then $\sum_i |b_i| > 0$ and $y^tA_iy = 0$ for all $y \in \R^m$ with $\sum x_iA_i y = 0$. Then we see that 
    \begin{align*}
        \sum x_iA_i \cdot Y = \sum x_ib_i \leq 0
    \end{align*}
    But since $\sum x_iA_i, Y \succeq 0$ the above must actually equal 0. Now we factor $Y = \sum \lambda_i v_iv_i^T$ via the spectral theorem. Then we see that,
    \begin{align*}
        \sum x_iA_i \cdot Y = \sum x_iA_i \cdot \sum \lambda_i v_iv_i^T = \sum \lambda_i \sum x_iA_i \cdot v_iv_i^T = 0
    \end{align*}
    Since a sum of positive numbers is zero iff each term is 0, we have that
    \begin{align*}
       \sum x_iA_i \cdot v_iv_i^T = 0
    \end{align*}

    We now show that if $A \succeq 0$ and $x \in \R^n$ is so that $x^TAx = 0$, then $x \in \ker A$. Write $A = \sum \lambda_i v_i$. We have:
    \begin{align*}
        0 = x^TAx = \sum \lambda_i x^Tv_iv_i^Tx = \sum \lambda_i (x^Tv_i)^2
    \end{align*}
    Thus $x^Tv_i = 0$ for all $i$. We see that $Ax = \sum_i \lambda_i v_i(v_i^Tx) = 0$, as desired. But then,
    \begin{align*}
        A_i \cdot Y = A_i \cdot \sum \lambda_i v_iv_i^T = \sum(A_i \cdot v_iv_i^T) = \sum(v_i^TA_i v_i) = 0
    \end{align*}
    A contradiction, since at least one of the $b_i$'s are nonzero. Thus the conditions cannot hold simultaneously. 

    Suppose that the second condition is false. Then if there is a matrix in $\mathcal A^*$ that is PSD, we are done. Suppose otherwise. Then we can find a matrix $M = \sum x_iA_i$ in $\widehat{\mathcal A^*}$ that is PSD. Since not all of the $b_i$ are 0, we need to show that $y^TA_iy = 0$ for all $y \in \ker M$. We make the following few simplifications. Replace $\mathcal A$ to be defined by new $A_i$'s with equations $A_i \cdot X = b_i$ for $A_i$ orthogonal (first give $\mathcal A$ a linearly independent basis then make that basis orthogonal by Gram Schmidt). Now let $y \in \ker M$. Since $\widehat{\mathcal A^*} = \mathrm{lin}(\mathcal A)^\perp$, we have that $\widehat{\mathcal A^*} \oplus \mathrm{lin}(\mathcal A)^\perp
\end{document}